% GENERATED FILE. Edit canonical lessons, then run npm run generate:books.
% canonical-source-hash: fnv1a64:fa28b3948e5613cc
% canonical-lessons: GU-C480-khas, GU-C480-jhakham, GU-C480-paru, GU-C480-sadvu, GU-C480-upasvu

\chapter{Trouble with the Skin}
\label{ch:gu-trouble-with-the-skin}
\hlchaptermodality{480}

\begin{chapteropening}
\textbf{By the end of this chapter:} \emph{I can describe a skin problem.}

Complete the last lesson of chapter 480: I can describe a skin problem.
\end{chapteropening}

\section[\texorpdfstring{\gu{ખસ} (khas)}{ખસ (khas)}]{\gu{ખસ} (khas) — an itchy rash; scabies}

\label{lesson:GU-C480-khas}



\begin{quote}
Before the new one: say the Gujarati for a lump, then the Gujarati for bedridden.
\end{quote}

\subsection*{You'll want to know: \gu{ખસ}}
\textbf{\gu{ખસ}} — \emph{khas} — \textquotedblleft{}an itchy rash; scabies\textquotedblright{}.

\begin{quote}
\gu{ગંદા} \gu{પાણીથી} \gu{બાળકને} \gu{ખસ} \gu{થઈ}.
\end{quote}

\begin{quote}
\emph{gandā pāṇīthī bāḷakne khas thaī.}
\end{quote}

\begin{quote}
The child got an itchy rash from dirty water.
\end{quote}

A feminine noun.

\begin{cousinweb}
\textbf{\gu{ખસ}} (\emph{khas}) goes back to Sanskrit \emph{khasa}, \textquotedblleft{}itch, scab\textquotedblright{}, which CDIAL thinks probably came from Dravidian; CDIAL lists the Gujarati word (CDIAL 3854).
\end{cousinweb}

\begin{grammarlens}[title={What you've built}]
One more word for describing what you lived and felt.
\end{grammarlens}

\subsection*{Your turn}
\begin{itemize}
\raggedright
  \item \emph{Say it:} \emph{khas}
  \item \emph{Say it:} \emph{khas}, once more
  \item \emph{Say it:} \emph{gandā pāṇīthī bāḷakne khas thaī.} — The child got an itchy rash from dirty water.
  \item \emph{Recall:} say the Gujarati for a lump, then the Gujarati for bedridden, then say \emph{khas} again
\end{itemize}

\begin{tcolorbox}[breakable,colback=teal!4,colframe=teal!35!black,title={Before you move on}]
What does \gu{ખસ} mean? (\textquotedblleft{}An itchy rash; scabies\textquotedblright{}.) Say it once more.
\end{tcolorbox}
\section[\texorpdfstring{\gu{ઝખમ} (jhakham)}{ઝખમ (jhakham)}]{\gu{ઝખમ} (jhakham) — a gash; a deep wound}

\label{lesson:GU-C480-jhakham}



\begin{quote}
Before the new one: say the Gujarati for bedridden, then the Gujarati for an itchy rash.
\end{quote}

\subsection*{You'll want to know: \gu{ઝખમ}}
\textbf{\gu{ઝખમ}} — \emph{jhakham} — \textquotedblleft{}a gash; a deep wound\textquotedblright{}.

\begin{quote}
\gu{છરીથી} \gu{તેના} \gu{હાથમાં} \gu{મોટો} \gu{ઝખમ} \gu{થયો}.
\end{quote}

\begin{quote}
\emph{charīthī tenā hāthmā̃ moṭo jhakham thayo.}
\end{quote}

\begin{quote}
He got a big gash on his hand from the knife.
\end{quote}

A masculine noun.

\begin{cousinweb}
\textbf{\gu{ઝખમ}} (\emph{jhakham}) is usually traced to Persian \emph{zaxm}, \textquotedblleft{}a wound\textquotedblright{}.
\end{cousinweb}

\begin{grammarlens}[title={What you've built}]
One more word for describing what you lived and felt.
\end{grammarlens}

\subsection*{Your turn}
\begin{itemize}
\raggedright
  \item \emph{Say it:} \emph{jhakham}
  \item \emph{Say it:} \emph{jhakham}, once more
  \item \emph{Say it:} \emph{charīthī tenā hāthmā̃ moṭo jhakham thayo.} — He got a big gash on his hand from the knife.
  \item \emph{Recall:} say the Gujarati for bedridden, then the Gujarati for an itchy rash, then say \emph{jhakham} again
\end{itemize}

\begin{tcolorbox}[breakable,colback=teal!4,colframe=teal!35!black,title={Before you move on}]
What does \gu{ઝખમ} mean? (\textquotedblleft{}A gash; a deep wound\textquotedblright{}.) Say it once more.
\end{tcolorbox}
\section[\texorpdfstring{\gu{પરુ} (paru)}{પરુ (paru)}]{\gu{પરુ} (paru) — pus}

\label{lesson:GU-C480-paru}



\begin{quote}
Before the new one: say the Gujarati for an itchy rash, then the Gujarati for a gash.
\end{quote}

\subsection*{You'll want to know: \gu{પરુ}}
\textbf{\gu{પરુ}} — \emph{paru} — \textquotedblleft{}pus\textquotedblright{}.

\begin{quote}
\gu{ઘામાં} \gu{પરુ} \gu{થયું} \gu{છે}, \gu{તેથી} \gu{દાક્તર} \gu{પાસે} \gu{જાઓ}.
\end{quote}

\begin{quote}
\emph{ghāmā̃ paru thayũ che, tethī dāktar pāse jāo.}
\end{quote}

\begin{quote}
The wound has pus in it, so go to the doctor.
\end{quote}

A neuter noun.

\begin{cousinweb}
\textbf{\gu{પરુ}} (\emph{paru}) goes back to a reconstructed \emph{*paripūya-}, \textquotedblleft{}purulent\textquotedblright{}; CDIAL lists the Gujarati word (CDIAL 7846).
\end{cousinweb}

\begin{grammarlens}[title={What you've built}]
One more word for describing what you lived and felt.
\end{grammarlens}

\subsection*{Your turn}
\begin{itemize}
\raggedright
  \item \emph{Say it:} \emph{paru}
  \item \emph{Say it:} \emph{paru}, once more
  \item \emph{Say it:} \emph{ghāmā̃ paru thayũ che, tethī dāktar pāse jāo.} — The wound has pus in it, so go to the doctor.
  \item \emph{Recall:} say the Gujarati for an itchy rash, then the Gujarati for a gash, then say \emph{paru} again
\end{itemize}

\begin{tcolorbox}[breakable,colback=teal!4,colframe=teal!35!black,title={Before you move on}]
What does \gu{પરુ} mean? (\textquotedblleft{}Pus\textquotedblright{}.) Say it once more.
\end{tcolorbox}
\section[\texorpdfstring{\gu{સડવું} (saḍvũ)}{સડવું (saḍvũ)}]{\gu{સડવું} (saḍvũ) — to rot; to decay}

\label{lesson:GU-C480-sadvu}



\begin{quote}
Before the new one: say the Gujarati for a gash, then the Gujarati for pus.
\end{quote}

\subsection*{You'll want to know: \gu{સડવું}}
\textbf{\gu{સડવું}} — \emph{saḍvũ} — \textquotedblleft{}to rot; to decay\textquotedblright{}.

\begin{quote}
\gu{મીઠાઈથી} \gu{દાંત} \gu{સડે} \gu{છે}.
\end{quote}

\begin{quote}
\emph{mīṭhāīthī dā̃t saḍe che.}
\end{quote}

\begin{quote}
Sweets make your teeth decay.
\end{quote}

Past (it …): \gu{તે} \gu{દાંત} \gu{સડી} \gu{ગયો}, \emph{te dā̃t saḍī gayo}.

\begin{cousinweb}
\textbf{\gu{સડવું}} (\emph{saḍvũ}) is inherited from Sanskrit \emph{śaṭati}, \textquotedblleft{}disintegrates\textquotedblright{}; CDIAL lists the Gujarati verb, \textquotedblleft{}to rot\textquotedblright{} (CDIAL 12268).
\end{cousinweb}

\begin{grammarlens}[title={What you've built}]
One more word for describing what you lived and felt.
\end{grammarlens}

\subsection*{Your turn}
\begin{itemize}
\raggedright
  \item \emph{Say it:} \emph{saḍvũ}
  \item \emph{Say it:} \emph{saḍvũ}, once more
  \item \emph{Say it:} \emph{mīṭhāīthī dā̃t saḍe che.} — Sweets make your teeth decay.
  \item \emph{Recall:} say the Gujarati for a gash, then the Gujarati for pus, then say \emph{saḍvũ} again
\end{itemize}

\begin{tcolorbox}[breakable,colback=teal!4,colframe=teal!35!black,title={Before you move on}]
What does \gu{સડવું} mean? (\textquotedblleft{}To rot; to decay\textquotedblright{}.) Say it once more.
\end{tcolorbox}
\section[\texorpdfstring{\gu{ઊપસવું} (ūpasvũ)}{ઊપસવું (ūpasvũ)}]{\gu{ઊપસવું} (ūpasvũ) — to puff up; to rise (of the skin)}

\label{lesson:GU-C480-upasvu}



\begin{quote}
Before the new one: say the Gujarati for pus, then the Gujarati for to rot.
\end{quote}

\subsection*{You'll want to know: \gu{ઊપસવું}}
\textbf{\gu{ઊપસવું}} — \emph{ūpasvũ} — \textquotedblleft{}to puff up; to rise (of the skin)\textquotedblright{}.

\begin{quote}
\gu{મચ્છર} \gu{કરડ્યા} \gu{પછી} \gu{ચામડી} \gu{ઊપસી} \gu{આવી}.
\end{quote}

\begin{quote}
\emph{macchar karaḍyā pachī cāmḍī ūpsī āvī.}
\end{quote}

\begin{quote}
After the mosquito bite the skin puffed up.
\end{quote}

Past (it …): \gu{તે} \gu{ફોલ્લો} \gu{ઊપસી} \gu{આવ્યો}, \emph{te phollo ūpsī āvyo}.

\begin{cousinweb}
\textbf{\gu{ઊપસવું}} (\emph{ūpasvũ}) goes back to a reconstructed \emph{*utpruṣati}, \textquotedblleft{}bubbles up\textquotedblright{}; CDIAL lists the Gujarati verb, \textquotedblleft{}to rise, of a blister or bubble\textquotedblright{} (CDIAL 1832).
\end{cousinweb}

\begin{grammarlens}[title={What you've built}]
One more word for describing what you lived and felt.
\end{grammarlens}

\subsection*{Your turn}
\begin{itemize}
\raggedright
  \item \emph{Say it:} \emph{ūpasvũ}
  \item \emph{Say it:} \emph{ūpasvũ}, once more
  \item \emph{Say it:} \emph{macchar karaḍyā pachī cāmḍī ūpsī āvī.} — After the mosquito bite the skin puffed up.
  \item \emph{Recall:} say the Gujarati for pus, then the Gujarati for to rot, then say \emph{ūpasvũ} again
\end{itemize}

\begin{tcolorbox}[breakable,colback=teal!4,colframe=teal!35!black,title={Before you move on}]
What does \gu{ઊપસવું} mean? (\textquotedblleft{}To puff up; to rise (of the skin)\textquotedblright{}.) Say it once more.
\end{tcolorbox}
